In this document, I will be outlining the scope of the computable semantic model by using a real example to motivate its features. The objective is to establish designs, reason about them and then use that to understand the scope of what the engine must be able to do. By doing this, the structure that needs to be implemented will reveal itself and there will principles will naturally arise that guide the implementation.

I will define basic computable semantic primitives and build the computable semantic model using that. The invariants specified by the CSP's will be enforced by reasoning about the design. For example, by going through this process, it will become clear that the validation mechanism will need to be able to find paths between nodes. In the same way, as I reason about the design I establish, I will make a list of features that I have to support.

This structured way of building this library will be more effective because it will force me to refine my thoughts and establish a structure that I will implement. It will also unambiguously define the scope of what is currently supported and how it can be expanded to support more features.